%% ============================================================
%%  TRA-SAE — Mixed-domain scientific reasoning with a 4B LLM
%%  Elsevier elsarticle format (numbered references)
%%  Build (Overleaf or local with elsarticle installed):
%%    pdflatex TRA_SAE_paper ; bibtex TRA_SAE_paper ; pdflatex ; pdflatex
%% ============================================================
\documentclass[preprint,12pt]{elsarticle}

\usepackage{amssymb}
\usepackage{amsmath}
\usepackage{booktabs}
\usepackage{multirow}
\usepackage{graphicx}
\usepackage{url}

%% ---- Convenience macros ----
\newcommand{\modelname}{\textsc{TRA-SAE}}
\newcommand{\basemodelname}{Qwen3.5-4B}
\newcommand{\dataset}{EXACT~2026}

\journal{Neurocomputing}

\begin{document}

\begin{frontmatter}

\title{\modelname{}: A Reproducible Low-Cost Curriculum and
       Reward-Shaping Study for Mixed-Domain Scientific Reasoning
       with a 4B Language Model}

\author[inst1]{Anonymous Author(s)}
\affiliation[inst1]{organization={Affiliation withheld for review},
            addressline={},
            city={},
            postcode={},
            state={},
            country={}}

\begin{abstract}
Scientific question answering couples numerical physical reasoning with
symbolic logical inference, and resource-constrained settings frequently
preclude large proprietary models or extensive reinforcement learning.
This paper reports \modelname{}, a low-cost curriculum training pipeline
for the \dataset{} mixed-domain scientific question answering task built
on a 4B-parameter open model. The pipeline combines general supervised
fine-tuning (SFT), a logic-focused curriculum adaptation stage, and Group
Relative Policy Optimization (GRPO) with lightweight format, correctness,
and unit-consistency reward terms; total training cost is \$18 and the full
pipeline including all evaluations is below \$64 on a single
A100-SXM4-80GB. On a 217-sample validation split, the strongest single
configuration raises overall accuracy from $36.41\%$ (zero-shot) to
$54.84\%$ ($69.50\%$ physics, $27.63\%$ logic); the largest single-stage
contribution is $+12.44$ percentage points (pp) from the first SFT stage.
McNemar tests confirm that the zero-shot-to-SFT and zero-shot-to-GRPO
transitions are statistically significant ($p<10^{-4}$ and $p<10^{-7}$),
whereas the GRPO-over-SFT margin is not ($p=0.21$); a three-seed evaluation
yields $52.69\pm1.92\%$ overall accuracy. Under a fair multi-prompt
re-evaluation, five 7B-class zero-shot baselines reach $11.5$--$29.0\%$,
confirming that the low baseline scores are not an artifact of prompt
formatting. On the external MMLU physics benchmark the fine-tuned model
generalizes ($48.62\%$ versus $42.29\%$ zero-shot), and a tool-augmented
variant (Python evaluation, Z3) yields a marginal $+0.46$~pp change. A
structured error analysis indicates that $69.4\%$ of residual errors are
answer-format extraction mismatches and $22.4\%$ are logical-fallacy
errors, which implies that reported accuracies are conservative and that
formal logic remains the principal reasoning bottleneck at this scale.
All code, checkpoints, and evaluation scripts are released.
\end{abstract}

\begin{highlights}
\item A three-stage curriculum (SFT, logic SFT, GRPO) raises a 4B model
      from $36.41\%$ to $54.84\%$ on mixed-domain scientific QA for under
      \$64 of total compute.
\item The first (general) SFT stage accounts for the largest single-stage
      gain ($+12.44$~pp); the GRPO-over-SFT increment is not statistically
      significant under the validation-set size.
\item A fair multi-prompt re-evaluation shows that 7B-class baselines
      remain at $11.5$--$29.0\%$, ruling out prompt formatting as the cause
      of their low scores.
\item Fine-tuning on \dataset{} transfers to the external MMLU physics
      benchmark ($+6.33$~pp over zero-shot).
\item Self-consistency voting and Dual-LoRA expert routing do not improve
      accuracy at 4B scale; most residual errors are answer-format
      extraction mismatches.
\end{highlights}

\begin{keyword}
large language models \sep scientific reasoning \sep
group relative policy optimization \sep curriculum fine-tuning \sep
physics question answering \sep logical reasoning \sep
parameter-efficient fine-tuning
\end{keyword}

\end{frontmatter}

%% ============================================================
\section{Introduction}
\label{sec:intro}
%% ============================================================

Large language models (LLMs) achieve strong results on mathematical and
logical reasoning benchmarks~\cite{Wei2022,Wang2023sc,Kojima2022}.
Deploying capable models under resource constraints, where larger
proprietary systems are unavailable, requires the joint design of the
training data curriculum, the reward specification, and the
parameter-efficient adaptation strategy.

Scientific reasoning presents a combined challenge. Physics problems
require formula selection, symbolic manipulation, and unit-consistent
numerical computation, whereas logic problems require quantifier handling,
propositional inference, and the detection of common fallacies. The two
modalities favor different inductive biases, yet a single model is expected
to address both.

The \dataset{} task provides a controlled testbed with $1{,}945$ training
examples ($1{,}213$ physics, $732$ logic) and a held-out evaluation set,
with a reported public baseline of $38.71\%$. This study treats the task
as a proxy for mixed-domain scientific question answering and quantifies
how far a \basemodelname{} model can be advanced using publicly available
fine-tuning techniques under a fixed compute budget.

The contributions of this paper are the following.
\begin{itemize}
  \item A low-cost three-stage curriculum pipeline (general SFT
        $\rightarrow$ logic curriculum SFT $\rightarrow$ GRPO reward
        shaping) for mixed-domain scientific question answering with a 4B
        open model, with a full end-to-end cost below \$64 on a single
        A100-80GB.
  \item An attribution of accuracy gains across pipeline stages: the first
        SFT stage contributes the largest single-stage gain
        ($+12.44$~pp over zero-shot, $p<10^{-4}$), the logic curriculum
        stage adds $+3.22$~pp, and GRPO adds $+2.77$~pp; the latter two
        increments are not individually significant under the validation
        set size, and this uncertainty is reported explicitly.
  \item A fair multi-prompt baseline re-evaluation in which each 7B-class
        model is evaluated under its native answer format, establishing
        that the observed baseline scores reflect task difficulty rather
        than prompt-format non-compliance.
  \item An external-benchmark evaluation on MMLU physics that quantifies
        the transfer of in-domain fine-tuning to an unseen test
        distribution, and a tool-augmented evaluation that quantifies the
        effect of Python and Z3 assistance on the best checkpoint.
  \item A characterization of negative results (self-consistency voting and
        Dual-LoRA expert routing) and a structured error taxonomy that
        separates answer-format extraction mismatches from genuine
        reasoning errors.
\end{itemize}

The remainder of the paper is organized as follows.
Section~\ref{sec:related} reviews related work.
Section~\ref{sec:problem} formalizes the task.
Section~\ref{sec:method} describes the \modelname{} pipeline.
Section~\ref{sec:setup} presents the experimental setup.
Section~\ref{sec:results} reports the results and analysis.
Section~\ref{sec:discussion} discusses limitations.
Section~\ref{sec:conclusion} concludes.

%% ============================================================
\section{Related Work}
\label{sec:related}
%% ============================================================

\subsection{Chain-of-Thought and Self-Consistency}
Wei et al.~\cite{Wei2022} showed that step-by-step prompting improves LLM
performance on arithmetic and symbolic reasoning, and
Kojima et al.~\cite{Kojima2022} demonstrated zero-shot chain-of-thought
prompting. Wang et al.~\cite{Wang2023sc} proposed self-consistency
decoding, which samples multiple reasoning paths and aggregates the most
frequent answer. The present results indicate that self-consistency
provides no net benefit at 4B scale on this task
(Section~\ref{sec:results}).

\subsection{Mathematics and Science Language Models}
DeepSeekMath~\cite{Shao2024} combined large-scale mathematical
pre-training with GRPO. Qwen2.5-Math~\cite{Yang2024} and the broader Qwen
family~\cite{Qwen2025} extended this approach across model scales.
Llemma~\cite{Azerbayev2024} performed continued pre-training on scientific
and proof corpora. MetaMath~\cite{Yu2023} and MAmmoTH~\cite{Yue2023}
constructed augmented mathematical reasoning datasets. The present work
operates at a smaller scale (4B) in a two-domain setting (physics and
logic) and introduces a unit-consistency reward term.

\subsection{Neuro-Symbolic Integration}
Logic-LM~\cite{Pan2023} and LINC~\cite{Olausson2023} couple LLMs with
formal solvers (Z3, first-order provers) for logical reasoning, and
PAL~\cite{Gao2022} delegates arithmetic to a Python interpreter.
Math-Shepherd~\cite{Wang2024mathshepherd} and process-supervision
methods~\cite{Lightman2023} provide step-level reward labels. This study
integrates the Z3 SMT solver into the correctness reward and into a
tool-augmented inference variant, while relying on outcome-level rewards to
avoid step-level annotation.

\subsection{Parameter-Efficient Fine-Tuning}
LoRA~\cite{Hu2022} decomposes weight updates into low-rank matrices, and
QLoRA~\cite{Dettmers2023} extends it to quantized backbones.
LoRAHub~\cite{Huang2023} composes multiple domain adapters. The Dual-LoRA
configuration evaluated here adapts this line of work to domain-specialized
routing.

\subsection{Reinforcement Learning for Language Models}
InstructGPT~\cite{Ouyang2022} established RLHF via PPO~\cite{Schulman2017}.
Direct Preference Optimization~\cite{Rafailov2023} removes the explicit
reward model. GRPO~\cite{Shao2024} computes relative rewards within a
generation group and omits a separate value network, which makes
single-GPU training of 4B-class models practical. Curriculum
learning~\cite{Bengio2009} motivates the staged data ordering used in the
present pipeline.

%% ============================================================
\section{Problem Formulation}
\label{sec:problem}
%% ============================================================

Let $\mathcal{D}=\mathcal{D}_{\text{phys}}\cup\mathcal{D}_{\text{logic}}$
denote the \dataset{} dataset. Each sample
$(q_i,a_i^{*},s_i)\in\mathcal{D}$ comprises a question $q_i$, a
ground-truth answer $a_i^{*}$, and a domain label
$s_i\in\{\text{physics},\text{logic}\}$. A model $M$ with parameters
$\theta$ maps a question to a free-text response
$\hat{y}_i=M_\theta(q_i)$. A structured extraction operator recovers the
candidate answer,
\begin{equation}
  a_i=\texttt{extract}(\hat{y}_i),
\end{equation}
and accuracy is
\begin{equation}
  \mathrm{Acc}=\frac{1}{|\mathcal{D}_{\text{val}}|}
  \sum_{i=1}^{|\mathcal{D}_{\text{val}}|}
  \mathbf{1}\!\left[\mathrm{verify}(a_i,a_i^{*},s_i,q_i)\right],
\end{equation}
where $\mathrm{verify}(\cdot)$ applies Z3-based symbolic checking for logic
samples and SI-unit-normalized numerical comparison for physics samples.

%% ============================================================
\section{Methodology}
\label{sec:method}
%% ============================================================

\subsection{Structured Response Format}
All models, base and fine-tuned, are instructed to produce a three-tag
response,
\begin{equation}
  \langle\texttt{reasoning}\rangle\ldots\langle/\texttt{reasoning}\rangle\;
  \langle\texttt{answer}\rangle\ldots\langle/\texttt{answer}\rangle\;
  \langle\texttt{explanation}\rangle\ldots\langle/\texttt{explanation}\rangle,
\end{equation}
which supports reliable answer extraction and enables a format-based reward
term.

\subsection{Phase 1: Supervised Fine-Tuning}
\basemodelname{} is fine-tuned on the full training set ($1{,}945$ samples)
with a cross-entropy loss on the target response. The configuration is
3~epochs, per-device batch size 4 with gradient accumulation 4 (effective
batch 16), learning rate $2\times10^{-4}$, LoRA $r=32$, $\alpha=64$, dropout
$0.05$, applied to all seven projection matrices. Training completes in
$58.8$~min on one A100-80GB and converges to a training loss of $0.3085$.

\subsection{Phase 1.5: Logic Curriculum Fine-Tuning}
A second SFT stage is restricted to the logic subset ($732$ samples) with a
reduced learning rate of $5\times10^{-5}$ for two epochs. This curriculum
stage~\cite{Bengio2009} allocates additional capacity to formal reasoning,
which is under-represented ($37.6\%$ of the training data). Training
completes in $15.7$~min and converges to a loss of $0.1326$.

\subsection{Phase 2: GRPO with Reward Shaping}
\paragraph{Reward function}
For a completion $y$ with ground truth $a^{*}$ and domain $s$, the reward is
\begin{equation}
  R(y,a^{*},s)=
    w_f\,r_{\text{format}}(y)
   +w_c\,r_{\text{correct}}(y,a^{*},s)
   +w_u\,r_{\text{unit}}(y,a^{*})
   +\lambda\,\mathbf{1}[\,|y_{\text{reason}}|>800\,],
  \label{eq:reward}
\end{equation}
with $w_f=0.30$, $w_c=0.60$, $w_u=0.10$, and $\lambda=-0.10$, where
$|y_{\text{reason}}|$ is the reasoning token count. The format term awards
$0.10$ per present tag (maximum $0.30$). The correctness term is binary and
applies a $\le 5\%$ relative tolerance after SI-unit normalization for
physics and normalized string or Z3-verified matching for logic. The
unit-consistency term equals $1$ when the extracted answer carries a unit
dimensionally compatible with the ground-truth unit, and $0$ otherwise.

\paragraph{GRPO training}
GRPO~\cite{Shao2024} is applied from the Phase~1.5 checkpoint with $K=4$
completions per prompt. The group-normalized advantage is
\begin{equation}
  \hat{r}_k=\frac{R_k-\bar{R}}{\sigma_R+\epsilon},
\end{equation}
with KL coefficient $\beta=0.04$, learning rate $10^{-6}$, and $250$ steps.
Training completes in $79.8$~min and reaches a loss of $0.0295$.

\subsection{Dual-LoRA Expert Routing}
As an additional configuration, two LoRA adapters are maintained, one
physics-specialized and one logic-specialized, and each question is routed
by a TF-IDF and logistic-regression classifier ($5{,}000$ n-gram features
fitted on the $1{,}945$ training samples). When the router confidence
exceeds $0.80$ the corresponding adapter is selected; otherwise a keyword
heuristic is used. Self-consistency voting with $N=5$ samples is applied as
a post-processing step in the full-agent configuration.

%% ============================================================
\section{Experimental Setup}
\label{sec:setup}
%% ============================================================

\subsection{Dataset}
Table~\ref{tab:dataset} summarizes the data split. The raw sources are a
logic corpus and a physics problem set, combined and divided by a $90/10$
stratified train/validation split.

\begin{table}[t]
\centering
\caption{Dataset statistics for \dataset{}.}
\label{tab:dataset}
\begin{tabular}{lrrr}
\toprule
Split & Total & Physics & Logic \\
\midrule
Train & $1{,}945$ & $1{,}213$ & $732$ \\
Validation & $217$ & $141$ & $76$ \\
\midrule
Total & $2{,}162$ & $1{,}354$ & $808$ \\
\bottomrule
\end{tabular}
\end{table}

\subsection{Baselines}
Six zero-shot baselines are evaluated on the $217$ validation samples:
Qwen2.5-Math-7B-Instruct~\cite{Yang2024}, Llemma-7B~\cite{Azerbayev2024},
Mistral-7B-Instruct-v0.3, Qwen2-Math-7B-Instruct, DeepSeek-Math-7B-Instruct,
and Gemini-2.5-flash-lite. The first five are open-weight models executed
locally; the last is a proprietary model accessed through an API and serves
as a reference point. All baselines receive the same system prompt and the
same verification protocol.

\subsection{Evaluation Metric}
The primary metric is overall accuracy; physics and logic sub-accuracies are
reported separately. Statistical significance between configuration pairs is
assessed by McNemar's test (two-sided, continuity-corrected), and Wilson
score intervals are reported for overall accuracy.

\subsection{Implementation Details}
Table~\ref{tab:hyperparams} lists the principal hyper-parameters. Decoding
uses greedy generation with a maximum of $1{,}024$ new tokens for all
configurations and baselines, which ensures that long chains of reasoning
reach the final answer before truncation.

\begin{table}[t]
\centering
\caption{Principal hyper-parameters.}
\label{tab:hyperparams}
\begin{tabular}{ll}
\toprule
Parameter & Value \\
\midrule
Base model & \basemodelname{} (BF16) \\
LoRA rank $r$ / $\alpha$ / dropout & $32$ / $64$ / $0.05$ \\
SFT learning rate / epochs & $2\times10^{-4}$ / $3$ \\
GRPO learning rate / steps & $1\times10^{-6}$ / $250$ \\
GRPO KL coefficient $\beta$ & $0.04$ \\
GRPO completions per prompt $K$ & $4$ \\
Reward weights $(w_f,w_c,w_u)$ & $(0.30,0.60,0.10)$ \\
Length penalty $\lambda$ & $-0.10$ if ${>}800$ tokens \\
Max new tokens & $1{,}024$ \\
GPU & A100-SXM4-80GB \\
\bottomrule
\end{tabular}
\end{table}

%% ============================================================
\section{Results and Analysis}
\label{sec:results}
%% ============================================================

\subsection{Main Results}
Table~\ref{tab:main} reports the full comparison. The strongest
configuration, cfg3 (SFT~$+$~logic~SFT~$+$~GRPO-mixed), reaches $54.84\%$
overall accuracy, $18.43$~pp above the \basemodelname{} zero-shot baseline
and above all six zero-shot baselines. The configurations are evaluated
with up to two answer-regeneration attempts, which is the agent evaluation
protocol used throughout the main comparison.

\begin{table}[t]
\centering
\caption{Main results on the \dataset{} validation set ($217$ samples:
         $141$ physics, $76$ logic). The double dagger ($\ddagger$) marks
         McNemar $p<0.01$ versus cfg0 for the two transitions reported in
         Table~\ref{tab:mcnemar} (cfg0$\rightarrow$cfg1,
         cfg0$\rightarrow$cfg3).}
\label{tab:main}
\begin{tabular}{llrrr}
\toprule
& System & Overall & Physics & Logic \\
\midrule
\multicolumn{5}{l}{\textit{Zero-shot baselines}} \\
B1 & \basemodelname{} (4B)            & $36.41\%$ & $45.39\%$ & $19.74\%$ \\
B2 & Qwen2.5-Math-7B-Instruct         & $28.57\%$ & $36.17\%$ & $14.47\%$ \\
B3 & Llemma-7B                        & $12.90\%$ & $\phantom{0}5.67\%$ & $26.32\%$ \\
B4 & Mistral-7B-Instruct-v0.3         & $\phantom{0}9.68\%$ & $\phantom{0}6.38\%$ & $15.79\%$ \\
B5 & Qwen2-Math-7B-Instruct           & $26.73\%$ & $33.33\%$ & $14.47\%$ \\
B6 & DeepSeek-Math-7B-Instruct        & $11.98\%$ & $\phantom{0}9.93\%$ & $15.79\%$ \\
B8 & Gemini-2.5-flash-lite (API)      & $24.88\%$ & $32.62\%$ & $10.53\%$ \\
\midrule
\multicolumn{5}{l}{\textit{Ours (fine-tuned \basemodelname{})}} \\
cfg0 & Zero-shot                      & $36.41\%$ & $45.39\%$ & $19.74\%$ \\
cfg1 & $+$ SFT Phase~1$^{\ddagger}$   & $48.85\%$ & $63.12\%$ & $22.37\%$ \\
cfg2 & $+$ Logic SFT                  & $52.07\%$ & $65.25\%$ & $27.63\%$ \\
cfg3 & $+$ GRPO-mixed$^{\ddagger}$    & $\mathbf{54.84\%}$ & $\mathbf{69.50\%}$ & $27.63\%$ \\
cfg4 & Dual-LoRA $+$ Router           & $50.69\%$ & $62.41\%$ & $28.95\%$ \\
cfg5 & Full Agent (SC $\times 5$)     & $52.53\%$ & $65.25\%$ & $28.95\%$ \\
\bottomrule
\end{tabular}
\end{table}

\subsection{Statistical Confidence and Significance}
Table~\ref{tab:ci} reports Wilson score intervals. The intervals for
cfg1--cfg5 overlap substantially, which indicates that the differences among
fine-tuned configurations are small relative to the validation-set size. The
non-overlap between the cfg0 and cfg1 intervals supports the conclusion that
the zero-shot-to-SFT gain is the robust effect.

\begin{table}[t]
\centering
\caption{Overall accuracy with $95\%$ Wilson score intervals ($n=217$).
         A representative subset of baselines is shown; the full set
         appears in Table~\ref{tab:main}.}
\label{tab:ci}
\begin{tabular}{llr}
\toprule
& System & Overall (\%) and $95\%$ CI \\
\midrule
B2 & Qwen2.5-Math-7B   & $28.57$\;\;[$23.0,34.9$] \\
B5 & Qwen2-Math-7B     & $26.73$\;\;[$21.3,33.0$] \\
B8 & Gemini-2.5-flash-lite & $24.88$\;\;[$19.6,31.0$] \\
B6 & DeepSeek-Math-7B  & $11.98$\;\;[$\phantom{0}8.3,17.0$] \\
\midrule
cfg0 & Zero-shot       & $36.41$\;\;[$30.3,43.0$] \\
cfg1 & $+$ SFT          & $48.85$\;\;[$42.3,55.5$] \\
cfg2 & $+$ Logic SFT    & $52.07$\;\;[$45.4,58.6$] \\
cfg3 & $+$ GRPO-mixed   & $54.84$\;\;[$48.2,61.3$] \\
cfg5 & Full Agent       & $52.53$\;\;[$45.9,59.1$] \\
\bottomrule
\end{tabular}
\end{table}

Table~\ref{tab:mcnemar} reports the pairwise McNemar tests computed on the
canonical per-sample predictions. The zero-shot-to-SFT transition
(cfg0~$\rightarrow$~cfg1) and the zero-shot-to-GRPO transition
(cfg0~$\rightarrow$~cfg3) are both significant. The logic-SFT increment
(cfg1~$\rightarrow$~cfg2, $p=0.10$) and the GRPO-over-SFT increment
(cfg2~$\rightarrow$~cfg3, $p=0.21$) are not significant at the
$0.05$ level; the self-consistency change
(cfg3~$\rightarrow$~cfg5) is a non-significant decrease.

\begin{table}[t]
\centering
\caption{Pairwise McNemar tests (two-sided, continuity-corrected,
         canonical per-sample predictions, $n=217$). $b$ is the count of
         samples where B is correct and A is wrong; $c$ is the reverse.}
\label{tab:mcnemar}
\begin{tabular}{lrrrrl}
\toprule
Pair (A~$\rightarrow$~B) & $\Delta$Acc & $b$ & $c$ & $\chi^2$ & $p$ \\
\midrule
cfg0~$\rightarrow$~cfg1 & $+12.44$ & $33$ & $\phantom{0}6$ & $17.33$ & $3.1\times10^{-5}$ \\
cfg0~$\rightarrow$~cfg3 & $+18.43$ & $44$ & $\phantom{0}4$ & $31.69$ & $<10^{-7}$ \\
cfg1~$\rightarrow$~cfg2 & $+3.22$  & $10$ & $\phantom{0}3$ & $2.77$  & $0.096$ \\
cfg2~$\rightarrow$~cfg3 & $+2.77$  & $11$ & $\phantom{0}5$ & $1.56$  & $0.211$ \\
cfg3~$\rightarrow$~cfg5 & $-2.31$  & $\phantom{0}5$ & $10$ & $1.07$ & $0.302$ \\
\bottomrule
\end{tabular}
\end{table}

\subsection{Multi-Seed Reliability}
Table~\ref{tab:multiseed} reports cfg3 re-evaluated with three generation
seeds. The overall standard deviation of $1.92$~pp is comparable to the
$+2.77$~pp GRPO-over-SFT increment, which reinforces the McNemar finding
that the GRPO margin should be interpreted with caution.

\begin{table}[t]
\centering
\caption{cfg3 multi-seed results (same checkpoint, three generation seeds).}
\label{tab:multiseed}
\begin{tabular}{lrrr}
\toprule
& Overall & Physics & Logic \\
\midrule
Seed 42   & $54.84\%$ & $69.50\%$ & $27.63\%$ \\
Seed 1337 & $51.15\%$ & $65.25\%$ & $25.00\%$ \\
Seed 2024 & $52.07\%$ & $65.96\%$ & $26.32\%$ \\
\midrule
Mean $\pm$ std & $52.69\pm1.92\%$ & $66.90\pm2.28\%$ & $26.32\pm1.32\%$ \\
\bottomrule
\end{tabular}
\end{table}

\subsection{Fair Multi-Prompt Baseline Re-evaluation}
\label{sec:fair}
The zero-shot baselines in Table~\ref{tab:main} use a single XML-format
prompt, which several models follow inconsistently. To separate format
compliance from reasoning capability, each baseline is re-evaluated under
three answer formats: XML tags, chain-of-thought with a boxed answer, and
chain-of-thought with a plain answer line. A unified extractor recovers the
answer in priority order, and the best score per model is reported in
Table~\ref{tab:fair}.

\begin{table}[t]
\centering
\caption{Fair multi-prompt re-evaluation. ``Best'' is the maximum overall
         accuracy across the XML, boxed, and plain answer formats; the
         selecting format is indicated ($n=217$).}
\label{tab:fair}
\begin{tabular}{llrr}
\toprule
Model & Best format & Best overall & (Physics / Logic) \\
\midrule
Qwen2-Math-7B-Instruct   & XML        & $29.03\%$ & $35.46$ / $17.11$ \\
Qwen2.5-Math-7B-Instruct & XML        & $24.88\%$ & $34.75$ / $\phantom{0}6.58$ \\
Mistral-7B-Instruct-v0.3 & plain CoT  & $16.59\%$ & $13.48$ / $22.37$ \\
DeepSeek-Math-7B-Instruct& boxed CoT  & $12.44\%$ & $14.89$ / $\phantom{0}7.89$ \\
Llemma-7B                & XML        & $11.52\%$ & $\phantom{0}4.96$ / $23.68$ \\
\bottomrule
\end{tabular}
\end{table}

Under the best-of-three protocol the strongest baseline reaches $29.03\%$,
and the native answer format improves only Mistral-7B (by $+6.91$~pp over
its XML score); the other models do not benefit. The low baseline scores
therefore reflect task difficulty rather than format non-compliance, and
the fine-tuned 4B model (cfg3, $54.84\%$) retains a substantial margin over
all 7B-class baselines under this fairer protocol.

\subsection{External Generalization: MMLU Physics}
\label{sec:mmlu}
To assess transfer beyond the in-domain split, cfg0, cfg2, and cfg3 are
evaluated on the MMLU~\cite{Hendrycks2021} high-school physics and college
physics subsets ($253$ four-option multiple-choice questions). No additional
fine-tuning is performed. Table~\ref{tab:mmlu} reports the results.

\begin{table}[t]
\centering
\caption{Zero-shot transfer to MMLU physics ($253$ questions:
         $151$ high-school, $102$ college). No MMLU-specific training.}
\label{tab:mmlu}
\begin{tabular}{lrrr}
\toprule
Configuration & Overall & High-school & College \\
\midrule
cfg0 (zero-shot)     & $42.29\%$ & $41.06\%$ & $44.12\%$ \\
cfg2 ($+$ logic SFT) & $47.04\%$ & $47.68\%$ & $46.08\%$ \\
cfg3 ($+$ GRPO)      & $\mathbf{48.62\%}$ & $46.36\%$ & $51.96\%$ \\
\bottomrule
\end{tabular}
\end{table}

The fine-tuned configurations outperform the zero-shot model on the external
benchmark by up to $+6.33$~pp, with monotonic improvement across the pipeline
stages. This indicates that the in-domain curriculum transfers to an unseen
physics test distribution rather than overfitting to the \dataset{} format.

\subsection{Tool-Augmented Reasoning}
\label{sec:tool}
Table~\ref{tab:tool} compares direct inference with a tool-augmented variant
of the best checkpoint, in which a Python evaluator handles physics
arithmetic and the Z3 solver handles logic inference; the language-model
answer is retained when the solver abstains. Both arms use single-pass
greedy decoding.

\begin{table}[t]
\centering
\caption{Direct versus tool-augmented inference on the best checkpoint
         ($n=217$, single-pass). Tool activations: Python evaluator on
         $62$ samples; the Z3 solver abstained on the remaining cases.}
\label{tab:tool}
\begin{tabular}{lrrr}
\toprule
Arm & Overall & Physics & Logic \\
\midrule
A: direct                 & $39.63\%$ & $49.65\%$ & $21.05\%$ \\
B: $+$ Python / Z3 tools  & $40.09\%$ & $51.06\%$ & $19.74\%$ \\
\midrule
$\Delta$ (B $-$ A)        & $+0.46$   & $+1.41$   & $-1.31$ \\
\bottomrule
\end{tabular}
\end{table}

Tool augmentation yields a marginal overall change ($+0.46$~pp), with a small
physics gain from the Python evaluator and a small logic decrease. The Z3
solver abstained on the natural-language logic items because their premises
were not reliably parseable into a formal specification, which limits the
benefit of symbolic tooling under the current parser.

\subsection{Reward Component Ablation}
Table~\ref{tab:reward} reports four reward variants trained for $150$ GRPO
steps from the Phase~1 checkpoint and evaluated under single-pass decoding,
which isolates the reward effect from the multi-attempt agent protocol used
in Table~\ref{tab:main}. The correctness-only variant (R4) attains the
highest accuracy at this budget, $1.84$~pp above the full mixture (R1).
Removing the format term (R3) reduces physics accuracy by $2.84$~pp.

\begin{table}[t]
\centering
\caption{Reward component ablation ($150$ GRPO steps from the Phase~1
         checkpoint, single-pass, $n=217$).}
\label{tab:reward}
\begin{tabular}{llrrr}
\toprule
& Reward configuration & Overall & Physics & Logic \\
\midrule
R1 & Full (format$+$correct$+$unit$+$length) & $39.63\%$ & $49.65\%$ & $21.05\%$ \\
R2 & No unit term                            & $39.17\%$ & $49.65\%$ & $19.74\%$ \\
R3 & No format term                          & $38.25\%$ & $46.81\%$ & $22.37\%$ \\
R4 & Correctness only                        & $\mathbf{41.47\%}$ & $\mathbf{51.06\%}$ & $23.68\%$ \\
\bottomrule
\end{tabular}
\end{table}

At a $150$-step budget the auxiliary unit-consistency and format terms do not
provide a net accuracy benefit relative to a correctness-only objective. A
plausible explanation is that the correctness reward supplies the most direct
training signal and that auxiliary terms add gradient variance that is not
amortized within $150$ steps; verifying whether the ordering reverses at
longer horizons was not feasible within the compute budget.

\subsection{Compute Cost}
Table~\ref{tab:compute} summarizes the compute profile. Phase~1 SFT provides
the highest accuracy per training dollar, and the full pipeline cost is below
one A100-day.

\begin{table}[t]
\centering
\caption{Compute profile on one A100-SXM4-80GB (\$3.67 per GPU-hour).}
\label{tab:compute}
\begin{tabular}{lrr}
\toprule
Stage & Time (min) & Cost \\
\midrule
SFT Phase 1            & $58.8$  & \$3.60 \\
Logic SFT (Phase 1.5)  & $15.7$  & \$0.96 \\
GRPO Phase 2 (mixed)   & $79.8$  & \$4.89 \\
GRPO Phase 2 (physics) & $76.6$  & \$4.69 \\
GRPO Phase 2 (logic)   & $62.9$  & \$3.85 \\
\midrule
Total training         & $293.8$ & \textbf{\$18.0} \\
\bottomrule
\end{tabular}
\end{table}

\subsection{Error Analysis}
\label{sec:error}
Table~\ref{tab:errors} reports a rule-based taxonomy of the $98$ cfg3 wrong
predictions. The largest category is answer-format extraction mismatch
($69.4\%$): the model produces a response whose format the automated verifier
cannot reconcile with the ground truth, for example a numerical value where
the ground truth is a qualitative label. The second category is
logical-fallacy and quantifier error ($22.4\%$), which is concentrated in the
logic subset.

\begin{table}[t]
\centering
\caption{Error taxonomy for cfg3 wrong predictions ($98$ of $217$).}
\label{tab:errors}
\begin{tabular}{llrr}
\toprule
ID & Error type & Count & \% \\
\midrule
E6 & Answer-format extraction mismatch   & $68$ & $69.4$ \\
E4 & Logical fallacy / quantifier error  & $22$ & $22.4$ \\
E0 & Unclassified                        & $\phantom{0}5$ & $\phantom{0}5.1$ \\
E1 & Unit or dimension error             & $\phantom{0}3$ & $\phantom{0}3.1$ \\
E2 & Wrong formula selection             & $\phantom{0}0$ & $\phantom{0}0.0$ \\
E3 & Arithmetic error                    & $\phantom{0}0$ & $\phantom{0}0.0$ \\
E5 & Question misinterpretation          & $\phantom{0}0$ & $\phantom{0}0.0$ \\
\bottomrule
\end{tabular}
\end{table}

Two implications follow. First, because most measured errors are extraction
mismatches rather than reasoning failures, the reported accuracies are
conservative lower bounds on the model's reasoning capability; the
\dataset{} answer space mixes numerical, categorical, and free-text targets,
which the automated verifier cannot always reconcile. Second, among genuine
reasoning errors the logical-fallacy category dominates, which is consistent
with the persistent physics--logic accuracy gap ($69.50\%$ versus $27.63\%$
for cfg3) and indicates that formal logic is the principal reasoning
bottleneck at this scale.

%% ============================================================
\section{Discussion}
\label{sec:discussion}
%% ============================================================

\subsection{Self-Consistency at 4B Scale}
Self-consistency is effective when errors are independent across
samples~\cite{Wang2023sc}. The present error profile indicates that a large
share of failures are systematic, either fixed answer-format mismatches or
repeated quantifier errors, which majority voting reinforces rather than
corrects. This accounts for the non-significant decrease from cfg3 to cfg5
($-2.31$~pp, $p=0.30$) despite five times the inference compute.

\subsection{Domain Routing}
The TF-IDF and logistic-regression router classifies clean physics and logic
items reliably, but border-zone questions, such as physics problems phrased
as logical conditionals, are routed to the wrong adapter. The Dual-LoRA
configuration (cfg4, $50.69\%$) does not exceed the single mixed adapter
(cfg3, $54.84\%$), which indicates that routing noise offsets the benefit of
domain specialization at this scale.

\subsection{Limitations}
\begin{itemize}
  \item \emph{Validation-set size.} With $217$ samples, Wilson intervals span
        roughly $\pm7$~pp, and differences below approximately $5$~pp are not
        statistically separable; the GRPO-over-SFT margin
        ($+2.77$~pp, $p=0.21$) falls in this regime.
  \item \emph{Budget-conditional reward finding.} The advantage of the
        correctness-only reward holds at a $150$-step budget and may not
        persist at longer horizons.
  \item \emph{Extraction-bounded measurement.} A majority of residual errors
        are answer-format extraction mismatches; the reported accuracies are
        lower bounds, and a manual annotation protocol is required to
        quantify the gap between measured and true reasoning accuracy.
  \item \emph{Symbolic tooling coverage.} The Z3 integration abstains on
        natural-language logic items that are not reliably parseable, which
        limits the benefit of the tool-augmented variant.
  \item \emph{Logic gap.} The physics--logic accuracy gap and the
        logical-fallacy error rate indicate that standard SFT and GRPO are
        insufficient for formal symbolic logic at 4B scale; step-level
        supervision~\cite{Lightman2023,Wang2024mathshepherd} or
        solver-guided training are the most direct avenues for improvement.
\end{itemize}

%% ============================================================
\section{Conclusion}
\label{sec:conclusion}
%% ============================================================

This paper presented \modelname{}, a low-cost curriculum training pipeline
for mixed-domain scientific reasoning that combines curriculum SFT, GRPO with
lightweight reward shaping, and an optional Dual-LoRA routing agent. The
strongest configuration reaches $54.84\%$ overall accuracy on the \dataset{}
validation set, an $18.43$~pp improvement over the \basemodelname{} zero-shot
baseline, using a 4B open model and under \$64 of total compute. The largest
single-stage contribution comes from the first (general) SFT stage
($+12.44$~pp), while the GRPO-over-SFT increment is not statistically
separable under the validation-set size. A fair multi-prompt re-evaluation establishes that
7B-class baselines remain at $11.5$--$29.0\%$, an external evaluation on MMLU
physics shows positive transfer ($+6.33$~pp over zero-shot), and a
tool-augmented variant produces only a marginal change. The error analysis
indicates that most residual errors are answer-format extraction mismatches
and that formal logic is the principal reasoning bottleneck. Future work
includes scaling to 7B and 14B checkpoints, process-level reward via
step-annotated data, cross-dataset evaluation on SciBench~\cite{Wang2024scibench}
and GPQA~\cite{Rein2023}, and tighter integration of solver-guided training
signals for the logic domain.

%% ============================================================
\section*{Data and Code Availability}
All training code, checkpoints, evaluation scripts, and result files are
released to support reproducibility.

%% ---- References ----
\bibliographystyle{elsarticle-num}
\bibliography{references}

\end{document}
